%% Motivic homotopy theory, as introduced by Morel and Voevodsky \cite{MorelVoevodsky}, has been a tremendously successful application of abstract homotopy theory to concrete questions of algebraic geometry and number theory. \todo{This sentence is stolen from Gheorghe--Wang--Xu. It should be rewritten.}

A striking surprise of modern computational stable homotopy theory is that the category $\SH(\C)$ of $\C$-motivic spectra, as introduced by Morel and Voevodsky \cite{MorelVoevodsky},
has found use not only in algebraic geometry, but also in the computation of stable homotopy groups of spheres, a question of purely topological origin.  In particular, the $p$-completed bi-graded homotopy groups of the unit in $\SH(\C)$ record---in a precise sense---the Adams--Novikov spectral sequence for the sphere spectrum, including all differentials and extensions.  The ultimate expression of this connection is the discovery that the \emph{cellular} subcategory of $p$-complete $\C$-motivic spectra is equivalent to a category of \emph{synthetic} Adams--Novikov spectral sequences, and in particular has a purely homotopy theoretic description without reference to algebraic geometry.

%% The relationship between $p$-complete, cellular $\C$-motivic spectra and Adams--Novikov spectral sequences is mediated by a class $\tau:\Ss^1 \to (\G_m)^{\wedge}_p$.  The cofiber of the $\tau$ map (often denoted $C\tau$) is an $\mathbb{E}_\infty$-algebra, with the property that cellular $C\tau$-modules record the purely algebraic derived category of $p$-complete $\BP_*\BP$-comodules.  On the other hand, the full subcategory of $\tau$-invertible, $p$-complete motivic $\C$-spectra is equivalent to the classical category of $p$-complete spectra.  One summarizes the situation by saying that $\SH(\C)^{cell,p}$ is a $\tau$-deformation, with generic fiber the category of $p$-complete spectra and special fiber the derived category of $p$-complete $\BP_*\BP$-comodules.

The first evidence that a large sector of the motivic category contains only topological information appears in the work of Thomason on algebraic K-theory and \'etale cohomology \cite{Thomason}. As a consequence of this work, Thomason proves that for well-behaved $\C$-schemes $X$, there is an equivalence $\mathrm{K}^{\mathrm{alg}} (X) [\beta^{-1}]^{\wedge} _{\ell} \simeq \mathrm{KU} (X(\C))^{\wedge} _{\ell}$, where $\beta$ is a certain Bott element. %% There is a precise sense in which inverting $\beta$ is equivalent to inverting $\tau$.
These ideas have been refined over the years, culminating in the following theorem.

\begin{thm}[\cite{VoeOpenProbs, DI, LevineComparison, LevineSlicemU, Ctau, ctauactalol, Pstragowski, cmmf}]\label{thm:C-main}\

  In $\SHC$
  there is a map $\tau : \Ss^1 \to (\G_m)_p$
  which enjoys the following properties:
  \begin{enumerate}
  \item[(1.\hspace{1em}]\hspace{-1.1em}$\mathrm{Generically\ Topological})$ The full subcategory of $\tau$-local objects in $\SHC^{cell}_p$ is equivalent to $\Sp_p$.
  \item[(2.\hspace{1em}]\hspace{-1.1em}$\mathrm{Algebraic\ Degeneration})$ The cofiber of $\tau$ (often denoted $C\tau$) is a commutative algebra, and the category of dualizable modules over $C\tau$ is equivalent to the category of $p$-completions of dualizable objects in $\QCoh(\mathcal{M}_{\mathrm{fg}})$, where $\Mfg$ is the moduli stack of formal groups.
  \item[(3.\hspace{1em}]\hspace{-1.1em}$\mathrm{Galois\ Reconstruction})$ There is a purely topological construction of $\SHC_p^{cell}$.
  \end{enumerate}
  
  We summarize this situation by saying that $\SH(\C)_p^{cell}$ is a $1$-parameter deformation of $p$-complete spectra with parameter $\tau$ and a purely algebraic special fiber.
\end{thm}

Although it does not appear in the statement, understanding the algebraic cobordism spectrum, $\MGL$, is key to proving this theorem. 
The importance of $\MGL$ to the study of cellular motivic spectra goes back to Voevodsky's work on the effective slice filtration \cite{VoeOpenProbs}.
In that work, Voevosky conjectured that the effective slice filtration of the unit over an algebraically closed field 
may be described in terms of the Adams-Novikov spectral sequence. 
%% may be described in terms of the $\mathrm{E}_2$-page of the Adams-Novikov spectral sequence and further suggested a relationship between the effective slice spectral sequence of the unit over an algebraically closed field and the Adams-Novikov spectral sequence.
Levine later proved Voevodsky's conjectures \cite{LevineComparison, LevineSlicemU}.
Voevodsky's notion of ``rigid homotopy groups'' also foreshadowed \emph{algebraic degeneration}.
In this framework he and Rezk predicted that the rigid Adams spectral sequence is equivalent to the algebraic Novikov spectral sequence \cite[p. 20-21]{VoeOpenProbs}, a result later proven in different language by Gheorghe--Wang--Xu \cite[Theorem 1.17]{ctauactalol}.
The fact that the category of $p$-complete $\C$-motivic spectra is \emph{generically topological} is strongly suggested by results of Dugger--Isaksen \cite[Section 2.6]{DI}, though to the best of our knowledge the full result did not appear in print until \cite{Pstragowski}.\todo{I can't find an earlier one at least, but that could easily just be me. Piotr does beat Tom's etale rigidity paper.}
The commutative algebra structure on $C\tau$ was first constructed by Gheorghe \cite{Ctau},
and the category of modules over $C\tau$ was then identified by Gheorghe--Wang--Xu in \cite{ctauactalol}.
Finally, % \emph{Galois reconstruction} was accomplished independently by 
Pstragowski \cite{Pstragowski} and Gheorghe--Isaksen--Krause--Ricka \cite{cmmf}
independently provided two different \emph{Galois reconstructions} of cellular $p$-complete $\C$-motivic spectra. %% moreover, \cite{Pstragowski} reproves (2) independently of \cite{Ctau} and \cite{ctauactalol}.

%% In SECTION, we show that there is a precise dictionary between Voevodsky's language of the effective slice filtration and rigid homotopy groups, and the language used in \Cref{thm:C-main}.

With the case of $\C$ resolved we raise the following natural question:

\begin{qst}\label{qst:main}
  To what extent can this theorem be extended to a general field $k$?
\end{qst}

As stated, this question is too imprecise to admit a definitive answer, so we begin by refining two points of ambiguity: the appropriate subcategory of $\SHk$ one should consider and the meaning of purely topological. In our study of this question we have found that the appropriate subcategory is the category of Artin--Tate motivic spectra (defined below). Notably, over $\R$, if one restricts to the further subcategory of cellular objects, it becomes significantly more difficult to construct a topological model.

\begin{dfn}
  The category of Artin--Tate motivic spectra over $k$ is the smallest stable full subcategory
  $\SHkat \subset \SHk$,
  closed under tensor products and colimits, that contains the motives of finite \'etale $k$-algebras, 
  $\P_k^1$ and $(\P_k^1)^{-1}$.
\end{dfn}

The phrase `purely topological' has a double meaning.
On the one hand it refers to the input to the construction;
it should not depend directly on the arithmetic of $k$, instead using only the absolute Galois group, $G$, together with the character $G \to \widehat{\Z}^\times$ induced by the maximal cyclotomic extension of $k$ \footnote{More specifically, the $p$-complete category should only depend on the $\Z_p^\times$ component of this map.}.
On the other hand it asks for a construction which uses only the ``standard machinery of homotopy theory.''

The authors are not the first to take up questions of this nature.
%% Indeed, beyond the work by many authors in the case $k = \C$ discussed above,
Positselski has studied the question of Galois reconstruction for the category $\DM(k;\F_\ell)^{\at}$ of mixed Artin-Tate motives with mod $\ell$ coefficients \cite{DMRecon,GalKosz}. With a few exceptions, he has shown that when $k$ is a finite, local or global field, then $\DM(k,\F_\ell)^{\at}$ may be viewed as a derived category of filtered discrete $G$-modules with restricted sub-quotients. %the restriction being given in terms of the cyclotomic character.

In a different direction, work of Bachmann, Elmanto and {\O}stv{\ae}r shows that, up to a completion, $\SH(S)$ is generically \'etale for a wide range of schemes $S$ \cite{TomEtaleI, TomEtaleII}. In particular, Bachmann--Elmanto--{\O}stv{\ae}r show that, after a suitable completion, \'etale localization corresponds to inverting $\tau$. Note that $\tau$ may not exist in the homotopy of the completed unit, so care must be taken to interpret this statement.
Furthermore, Bachmann showed that, again up to completion, the \'etale motivic category is equivalent to the category of hypercomplete sheaves of spectra on the small \'etale site.
Specializing to the case $S = \Spec k$, we find that, up to completion, the category of $\tau$-local objects in $\SH(k)$ admits a description in terms of Borel $G$-equivariant spectra.

Previously, work of Behrens and Shah \cite{BehrensShah} had taken up the question of when a suitable map $\tau$ exists over $\R$. 
Although there is no map $\tau:\Ss^1 \to (\G_m)^{\wedge}_2$ in $\SH(\R)$,
%% though such a class does exist after tensoring with the spectrum $\HFt$ representing mod $2$ motivic cohomology.
they prove that $\tau$ exists whenever a different class $\rho:\Ss^0 \to \G_m$ has been killed.
If only $\rho^2$ is killed, but not $\rho$ itself, then $\tau$ does not necessarily exist, but $\tau^2$ does.
Continuing in this way, they make sense of inverting $\tau$ in any $\rho$-complete situation.
%% but the resulting $\tau$-inverted objects are more closely related to Borel $C_2$-spectra than genuine $C_2$-spectra.
%% Furthermore, there is no way to construct $C\tau$ without first killing $\rho$. This makes it impossible to describe cellular $\R$-motivic spectra as a $\tau$-deformation.

%% \todo{What's below needs to be joined up with what's above.}

%% In recent years the category $\SHR$ of $\R$-motivic spectra has similarly been found to shed light on the genuine $C_2$-equivariant stable homotopy category.  However, the exact relationship between $\R$-motivic stems and any sort of equivariant Adams--Novikov spectral sequence has remained mysterious, and before the present work there has been no purely homotopy theoretic construction of the cellular subcategory of $\SH(\R)$.

In this paper we resolve \Cref{qst:main} in the case $k=\R$.
This begins with the observation that if one does not insist that the target of $\tau$ must be a completion of $\G_m$, then a suitable replacement can easily be constructed.
%% The first 
%% The main moral of this paper is that all of these subtleties disappear so long as one is willing to consider grading $\R$-motivic homotopy groups not just by the Picard elements $\Ss^1$ and $\mathbb{G}_m$, but also by a third Picard element $\Ss^{\C}$.  Here, $\Ss^{\C}$ is the suspension spectrum of the unreduced suspension of $\Spec(\C)$.

\begin{thm}\label{thm:main}
  In $\SHRat$
  there is an invertible object $Q$ and a map $\ta : \Ss^1 \to Q_2$
  which enjoys the following properties:%\footnote{As we will explain in our conventions below, we use $\SHRat$ to denote the category of modules over the $2$-completion of the unit in $\SHR^{\at}$.}
  \begin{enumerate}
  \item[(GT)] The full subcategory of $\ta$-local objects in $\SHR^{\at}_{i2}$ is equivalent to $\Sp_{C_2,i2}$.
  \item[(AD)] The cofiber of $\ta$ (denoted $C\ta$) is a commutative algebra, and the category of dualizable modules over $C\ta$ is equivalent to the category of dualizable objects in the derived category of Mackey-functor $\MU_*\MU$-comodules \footnote{See \Cref{sec:cta} for a precise definition of this category.}.
  \item[(GR)] There is a purely topological construction of $\SHRat$.
    In particular, we construct a commutative algebra $R_\bullet$ in filtered, $C_2$-equivariant spectra
    such that the category of filtered modules over $R_\bullet$ is equivalent to $\SHRat$.
    The commutative algebra $R_\bullet$ is the even slice--d\'ecalage of the $\MUR$-Adams tower for the sphere \footnote{See \Cref{sec:top} for precise definitions.}.    
  \end{enumerate}
  
  We summarize this by saying that $\SHR^{\at}_{i2}$ is a $1$-parameter deformation of $C_2$-equivariant stable homotopy theory with parameter $\ta$ and purely algebraic special fiber.
\end{thm}

The romanization of the Devanagari letter $\ta$ is `ta' \footnote{However, the reader who is not familiar with the Devanagari alphabet is advised that the pronunciation of $\ta$ in the majority of South Asian languages using Devanagari is given in IPA by /\textsubbridge{t}\textschwa/, which is closer to `tuh' in English than `ta.'}  and
our choice of this symbol will be explained in \Cref{rmk:ta-explanation}. As in the case over $\C$, we are also able to give an explicit formula for the homotopy groups of $C\ta$, though since we have not yet set up the appropriate notion of homotopy groups we will defer an explicit statement until later.
%% Note that the $p$-completions present in the theorem statement are necessary.
The subscript $i2$ appearing in the theorem statement refers to the category of modules over the $2$-completion of the unit. 
The reader who is worried by this departure from the norm may wish to read the section on notations and conventions before proceeding.
%% about the handling of $p$-completion may wish to read \Cref{app:completion} which sets up our general approach to completion issues before proceeding.

\ 

With the resolution of \Cref{qst:main} in the case of $\R$, the authors are hopeful that we might see this question resolved in its entirety in the near future.

%% \begin{cnstr}
%% There is a canonical class $\ta: \Ss^{\C} \to (\G_m)^{\wedge}_2$.  We will build $\ta$ in \Cref{sec:tau} via the explicit algebraic geometry of the circle $Q=\Spec \R[x,y]/(x^2+y^2-1)$.  A key point is that $\Sigma^{\infty} Q$ is itself a Picard element in $\SHgc$, satisfying a relation $\Ss^{\C} \otimes \Sigma^{\infty} Q \simeq \Ss^1 \otimes \G_m$.
%% \end{cnstr

%% \begin{rmk}
%%   The romanization of the Devanagari letter $\ta$ is `ta.' \footnote{However, the reader who does not know Hindi or another language which utilizes the the Devanagari alphabet should be advised that its pronunciation is closer to `tuh.' To be even more correct, the `t' sound should be dental and unaspirated.}  As we will explain in \Cref{sec:homology}, our choice of this character arises from the fact that there is a class $u:\Ss^1 \otimes (\Ss^{\C})^{-1} \to \HFt$ and a product relation $\ta u = \tau \in \pi_{***}^{\R}(\HFt).$    In other words, ta $\cdot$ u is tau.
	
%% It follows that inverting $\tau$ is the same as inverting both $\ta$ and $u$.  We have shown that inverting $\ta$ is equivalent to Betti realization, while inverting $u$ Borel completes genuine $C_2$ spectra.
